% =====================================================================
%  4.2 Den deriverte — fasit
% =====================================================================

\fk{Del 1 — Grenseverdi og definisjonen av den deriverte}
\fo{1.1} a) $f(5) = 25$ og $(5 + h)^2 = 25 + 10h + h^2$ \ b) $\dfrac{10h + h^2}{h} = 10 + h$ \ c) Når $h \to 0$: $f'(5) = 10$
\fo{1.2} $\dfrac{(2 + h)^2 + (2 + h) - 6}{h} = \dfrac{5h + h^2}{h} = 5 + h$. Når $h \to 0$: $f'(2) = 5$.
\fo{1.3} a) $\dfrac{5(4 + h) - 2 - 18}{h} = 5$. Når $h \to 0$: $f'(4) = 5$ \ b) $g'(1) = -3$. Hvert svar er stigningstallet til den rette linja.
\fo{1.4} a) $f'(x) = 2x + 6$ \ b) $f'(2) = 10$ og $f'(-1) = 4$
\fo{1.5} a) $\dfrac{s(3 + h) - s(3)}{h} = \dfrac{13h + 2h^2}{h} = 13 + 2h$. Når $h \to 0$: $s'(3) = 13$ \ b) Etter 3 sekunder er farten 13 m/s.
\fo{1.6} Feilen er overgangen til $\dfrac{16 + h^2 - 12 - 3h - 4}{h}$. Riktig er $(4 + h)^2 = 16 + 8h + h^2$, ikke $16 + h^2$. Da blir telleren $5h + h^2$, og kvotienten er $5 + h$. Når $h \to 0$: $f'(4) = 5$.
\fo{1.7} a) Vi kan ikke sette $h = 0$. Teller og nevner blir 0, og $\dfrac{0}{0}$ har ingen verdi. Det er ikke 0. Vi må forenkle og forkorte $h$ først, og så la $h \to 0$.
b) $f(a + h) - f(a)$ er endringen i $y$ når $x$ øker med $h$. Delt på $h$ får vi sekantens stigningstall. Grenseverdien når $h \to 0$ er tallet sekantene nærmer seg, altså $f'(a)$.
c) $f'(3)$ er stigningstallet til tangenten i $x = 3$. Det sier hvor bratt grafen er der: positivt tall betyr stigende, negativt betyr synkende.
\fo{1.8} a) Noen ganger. For $f(x) = x^2 - 6x$ er $f'(4) = 2 > 0$, men $f(4) = -8$. For $f(x) = x^2$ stemmer det i $x = 3$.
b) Alltid. $\dfrac{k(a + h) + m - ka - m}{h} = \dfrac{kh}{h} = k$ for alle $h \neq 0$.
c) Noen ganger. For $f(x) = x^2$ er $f'(0) = 0$ og $0$ er et bunnpunkt. For $f(x) = x^3$ er $\dfrac{h^3}{h} = h^2 \to 0$, så $f'(0) = 0$, men grafen stiger hele veien og har ikke topp- eller bunnpunkt.
\fo{1.9} a) $f'(1) = 8$ \ b) For eksempel $f(x) = 3x^2 + 2x$. Da er $f(1) = 5$ og $f(1 + h) = 5 + 8h + 3h^2$, så kvotienten er $8 + 3h$.
c) Ja, uendelig mange. Legger vi til en konstant $c$, forsvinner den i $f(1 + h) - f(1)$.
\fo{1.10} $\dfrac{f(2 + h) - f(2)}{h} = \dfrac{p(4h + h^2) - h}{h} = 4p + ph - 1$. Når $h \to 0$: $f'(2) = 4p - 1$. Fra $4p - 1 = 11$ får vi $p = 3$.
\fo{Sjekk 1} a) $\dfrac{4h + h^2}{h} = 4 + h$, så $f'(2) = 4$ \ b) $f(2) = 9$ er høyden til grafen. $f'(2) = 4$ er stigningstallet til tangenten. c) Teller og nevner blir 0, og $\dfrac{0}{0}$ har ingen verdi. Vi må forenkle og forkorte $h$ først.

\fk{Del 2 — Utlede derivasjonsregler}
\fo{2.1} a) $5x^4$ \ b) $9x^8$ \ c) $1$ \ d) $0$
\fo{2.2} a) $8 + 12h + 6h^2 + h^3$ \ b) $\dfrac{12h + 6h^2 + h^3}{h} = 12 + 6h + h^2$ \ c) Når $h \to 0$: $f'(2) = 12$
\fo{2.3} a) $10$ \ b) $3$ \ c) $2x = 8$ gir $x = 4$
\fo{2.4} $\dfrac{2(x + h)^2 + 5(x + h) - 2x^2 - 5x}{h} = \dfrac{4xh + 2h^2 + 5h}{h} = 4x + 2h + 5$. Når $h \to 0$: $g'(x) = 4x + 5$.
\fo{2.5} a) $s'(t) = 3t^2 + 2t$ (sumregelen, $(t^3)' = 3t^2$ og $(t^2)' = 2t$) \ b) $\dfrac{s(2 + h) - s(2)}{h} = \dfrac{16h + 7h^2 + h^3}{h} = 16 + 7h + h^2$. Når $h \to 0$: $s'(2) = 16$ m/s. Det stemmer med a): $3 \cdot 4 + 2 \cdot 2 = 16$.
\fo{2.6} Feilen er overgangen til $3x^2 + 3xh + h^2 + 4h$. Leddet $4h$ skal også deles på $h$ og blir $4$. Riktig: $3x^2 + 3xh + h^2 + 4$. Når $h \to 0$: $f'(x) = 3x^2 + 4$.
\fo{2.7} a) $x^7 + 7x^6h + 21x^5h^2 + 35x^4h^3 + 35x^3h^4 + 21x^2h^5 + 7xh^6 + h^7$ \ b) $7x^6$ \ c) Ja, det stemmer med $n = 7$. Men det er ett tilfelle til, ikke et bevis for alle $n$.
\fo{2.8} $\dfrac{k\,f(x + h) - k\,f(x)}{h} = k \cdot \dfrac{f(x + h) - f(x)}{h}$. Når $h \to 0$ går brøken mot $f'(x)$, så $(k \cdot f)' = k \cdot f'$.
\fo{2.9} a) Noen ganger. Med $f(x) = g(x) = x$ er $(x \cdot x)' = (x^2)' = 2x$, men $f' \cdot g' = 1 \cdot 1 = 1$. Med $f = g = 0$ er begge sider $0$.
b) Noen ganger. Med $f(x) = x^2 + 1$ og $g(x) = x^2$ er $f'(x) = g'(x) = 2x$, men $f(x) \neq g(x)$. Det stemmer når $f = g$. Funksjoner med samme deriverte kan skille seg med en konstant.
\fo{2.10} $f(x) = x^2 + 3x + 1$. Da er $f'(x) = 2x + 3$ og $f(0) = 1$. Uten kravet finnes uendelig mange: $x^2 + 3x + c$ for en hvilken som helst konstant $c$, fordi $(c)' = 0$.
\fo{Sjekk 2} a) $6x^5$ \ b) $2x$. Sumregelen $(f + g)' = f' + g'$ gir $(x^2)' + (9)' = 2x + 0$. Den stemmer fordi differansekvotienten til en sum er summen av differansekvotientene. \ c) $(2 + 1)^3 = 27$, men $2^3 + 1^3 = 9$
\fk{Ekstra utfordring}
\fo{U1} a) $f(2) = -1$, og $\dfrac{f(2 + h) - f(2)}{h} = \dfrac{h + h^2}{h} = 1 + h$. Når $h \to 0$: $f'(2) = 1$ \ b) Tangenten har stigningstall $1$, så $y = x + b$. Punktet $(2, -1)$ gir $-1 = 2 + b$, så $b = -3$ og $y = x - 3$.
\fo{U2} $\dfrac{(a + h)^2 - a^2}{h} = 2a + h$. Når $h \to 0$: $f'(a) = 2a$. Vi vil ha $2a = 8$, så $a = 4$. Punktet er $(4, 16)$.
\fo{U3} a) $6,001$ og $6,0$. Den første tilnærmingen er $6 + h$ med $h = 0,001$, så den er litt for stor. \ b) $x^3$: forover $12,006$, symmetrisk $12,0$. $1/x$: forover $-0,2499$, symmetrisk $-0,25$. Den symmetriske er nærmest. \ c) For $x^3$ i $x = 2$: forover gir $12,61$ og $12,0601$, symmetrisk gir $12,01$ og $12,0001$. Forover-feilen blir omtrent 10 ganger mindre når $h$ blir 10 ganger mindre. Den symmetriske feilen blir 100 ganger mindre.
\fo{U4} a) $(x^2 + 2xh + h^2)^2 = x^4 + 4x^3h + 6x^2h^2 + 4xh^3 + h^4$. Rad: $1\ 4\ 6\ 4\ 1$.
b) \[ \begin{aligned} \dfrac{(x + h)^4 - x^4}{h} &= \dfrac{4x^3h + 6x^2h^2 + 4xh^3 + h^4}{h}\\ &= 4x^3 + 6x^2h + 4xh^2 + h^3 \end{aligned} \]
Når $h \to 0$: $(x^4)' = 4x^3$.
c) Tallet er $n$ (nest etter 1 i rad $n$). Leddet $n\,x^{n-1}h$ er det eneste uten $h^2$ eller mer, så det blir igjen som $nx^{n-1}$ når vi deler på $h$ og lar $h \to 0$. Dette viser ideen. Et fullt bevis trenger binomialformelen.
\fo{U5} a)
\[ \dfrac{\dfrac{1}{x + h} - \dfrac{1}{x}}{h} = \dfrac{\dfrac{x - (x + h)}{x(x + h)}}{h} = \dfrac{-h}{x(x + h)\,h} = -\dfrac{1}{x(x + h)} \]
Når $h \to 0$: $f'(x) = -\dfrac{1}{x^2}$.
b) $f'(x) < 0$ for alle $x \neq 0$. Begge grenene av grafen synker.
\fo{U6} a) Utvid med $\sqrt{x + h} + \sqrt{x}$:
\[ \dfrac{(x + h) - x}{h\left(\sqrt{x + h} + \sqrt{x}\right)} = \dfrac{1}{\sqrt{x + h} + \sqrt{x}} \]
Når $h \to 0$: $f'(x) = \dfrac{1}{2\sqrt{x}}$.
b) Ja. Mønsteret gir $\frac{1}{2}x^{-\tfrac{1}{2}} = \dfrac{1}{2\sqrt{x}}$.
\fo{U7} a) $\dfrac{(x + h)^2 - (x - h)^2}{2h} = \dfrac{4xh}{2h} = 2x$ \ b) For $x^3$ er den symmetriske $3x^2 + h^2$ og den vanlige $3x^2 + 3xh + h^2$. Den symmetriske ligger nærmest for de fleste $x$. Flere funksjoner kan velges. c) Punktene ligger på hver sin side av $x$. Feilleddene med $h$ opphever hverandre, og vi sitter igjen med feil av størrelse $h^2$ i stedet for $h$. Sekanten ligner mer på tangenten.
